% Generated draft: requires mathematical and source review.
\documentclass[11pt]{article}
\usepackage[margin=25mm]{geometry}
\usepackage{fontspec}
\setmainfont{Latin Modern Roman}
\usepackage{amsmath,amssymb,mathtools,unicode-math}
\setmathfont{Latin Modern Math}
\usepackage{enumitem,needspace,xurl,hyperref}
\setlength{\parindent}{0pt}
\setlength{\parskip}{5pt}
\setlength{\emergencystretch}{3em}
\newcommand{\N}{\mathbb{N}}
\newcommand{\Z}{\mathbb{Z}}
\newcommand{\Q}{\mathbb{Q}}
\newcommand{\R}{\mathbb{R}}
\newcommand{\C}{\mathbb{C}}
\newcommand{\F}{\mathbb{F}}
\newcommand{\A}{\mathbb{A}}
\newcommand{\PP}{\mathbb{P}}
\newcommand{\E}{\mathbb{E}}
\newcommand{\eps}{\varepsilon}
\newcommand{\dd}{\,\mathrm{d}}
\begin{document}
\section*{UnsolvedMath ID: 30005278}
\textbf{Source problem code:} catalogue entry
\subsection{Definitions and mathematical statement}
For \(R\in\{\Z,\Q\}\), write \(\operatorname{Irr}_{R[x]}(h)\) if \(h\in R[x]\) is nonzero, is not a unit, and every factorization \(h=uv\) in \(R[x]\) has a unit factor. The units of \(\Z[x]\) are \(\pm1\), whereas the units of \(\Q[x]\) are all nonzero rational constants. Composition means \((f\circ g)(x)=f(g(x))\).

Define
\[
\mathcal G_{\leq2}=\{ax^2+bx+c:a,b,c\in\Z\},\qquad
\mathcal G_{1,2}=\{ax^2+bx+c:a,b,c\in\Z,\ (a,b)\ne(0,0)\}.
\]
Thus \(\mathcal G_{\leq2}\) includes constants and the zero polynomial, interpreting its degree bound by coefficient form; \(\mathcal G_{1,2}\) consists exactly of the polynomials of degree one or two.

The supplied reconstructed question asks whether there exists \(f\in\Z[x]\) of degree five whose composition with every integer polynomial of degree at most two is irreducible. With the coefficient-form interpretation above and a specified ring \(R[x]\), its literal formulation is
\[
\exists f\in\Z[x]\quad
\bigl(\deg f=5\ \land\
\forall g\in\mathcal G_{\leq2},\
\operatorname{Irr}_{R[x]}(f\circ g)\bigr).
\]
The original extracted question instead uses the term \emph{\(2\)-superirreducible}, without supplying its definition. If that term is intended to quantify only over nonconstant substitutions, the corresponding candidate formulation is
\[
\exists f\in\Z[x]\quad
\bigl(\deg f=5\ \land\
\forall g\in\mathcal G_{1,2},\
\operatorname{Irr}_{R[x]}(f\circ g)\bigr).
\]
These are conditional formulations, not an assertion that the source equates them. The intended substitution class, treatment of the zero polynomial, and irreducibility ring require clarification. In particular, constant substitutions produce constant compositions, whose irreducibility depends on the ring convention. The supplied evidence does not resolve these qualifications.
\subsection{Short English statement}
The question concerns a degree-five polynomial with integer coefficients whose compositions with all permitted integer polynomials of degree at most two remain irreducible. The supplied wording leaves unresolved whether substitutions must be nonconstant and whether irreducibility is taken in \(\Z[x]\) or \(\Q[x]\). A definitive statement of the intended \(2\)-superirreducibility problem requires these conventions.
\subsection{Sources}
\begin{itemize}
\item Supplied local UnsolvedMath registry record, permanent ID \(30005278\), source problem code \texttt{catalogue entry}, title \emph{Degree-Five Two-Superirreducible Polynomials}; collection: Oberwolfach Reports Open Problems.
\item Supplied report citation: \emph{Analytic Number Theory (2023)}, Oberwolfach Reports, DOI \texttt{10.4171/owr/2022/50}, pp.~2895--2960. \url{https://doi.org/10.4171/owr/2022/50}. The registry records proposed year \(2022\).
\item Supplied literature source: \emph{Oberwolfach Report 50}. \url{https://ems.press/content/serial-article-files/46986}.
\item Supplied literature source attributed to Bober, Du, Fretwell, Kopp, and Wooley. \url{https://doi.org/10.1016/j.indag.2024.08.005}. No article title is supplied.
\item Literature assessment dated exactly \texttt{2026-08-22}: the supplied curation labels the integral degree-five \(2\)-superirreducibility question open and unresolved, and describes recent work as addressing finite-field analogues rather than this integer-polynomial problem. This is supplied dated curation, not an independently verified assessment; its time-sensitive status and cited claims require independent verification.
\item The registry labels the reconstructed statement \texttt{corrected\_verified} and reports that it was checked against the official Oberwolfach report. This verification claim is supplied metadata and has not been independently checked here. The original extraction contains an incomplete attribution, ``a question of J.'', followed by unrelated problem material; the proposer cannot be identified from that fragment.
\end{itemize}
\end{document}
